\section{演练测试与正式测试结果}\label{sec:results}

\subsection{测试环境与客户端}

机器狗程序用 Python 实现，通过 HTTP$+$JSON 与官方模拟器通信。客户端遵守附件 2 的全部要求：以参赛队号作为 \texttt{robot\_id}，每个新动作使用新的 \texttt{request\_id}，\textbf{串行发送、逐次等待响应}，网络中断后重试时复用原请求内容与原 \texttt{request\_id}；结果未知的动作标记为 \texttt{uncertain} 并禁止发出不同的后续动作，以免重复计时或状态错位。客户端按 \texttt{/enter} 返回的剩余时长设置现实截止时间，并同步记录逐动作日志、配置哈希与定位证书，供事后核对。

\subsection{演练样本的选取规则}

演练测试不限次数，因此日志中混有不同算法版本的记录。为避免把旧版本的成绩算进当前方案，本文按三条规则筛选样本，规则在看到成绩之前就已固定：

\begin{enumerate}[label=(\arabic*),itemsep=1pt]
  \item 该局状态为 \texttt{completed} 且由机器狗\textbf{主动调用 \texttt{/exit}} 结束——中途失败、超时或被手工中止的记录一律剔除；
  \item 该局所用的策略源码哈希与当前提交版本一致（第三问 \texttt{d12fa3a5}，第四问 \texttt{3de0e9a0}）；
  \item 该局所用的配置与当前默认配置逐项一致。
\end{enumerate}

按此规则，第三问保留 $50$ 局、第四问保留 $70$ 局；另有 $25$ 局第三问记录因补测精化的参数与当前配置不同而被排除，排除清单与原因随支撑材料一并给出。\textbf{没有任何一局因为成绩不好而被剔除。}

\subsection{演练测试结果}

演练结束后模拟器界面会显示本局干扰源总数，逐局核对表明它与机器狗的成功清除数完全一致，故
\begin{equation}
  \text{被清除干扰源个数的比例}=\frac{N_c}{n}=100\%\quad\text{（全部 }120\text{ 局）} .
\end{equation}
汇总结果见\cref{tab:practice-summary}，按源数分层的结果见\cref{tab:practice-strat}。逐局明细（含每局的时间分项与动作数）见支撑材料 \texttt{q3\_runs.csv} 与 \texttt{q4\_runs.csv}。

\begin{table}[!htbp]
  \caption{官方模拟器演练测试汇总（时间单位：虚拟秒）}\label{tab:practice-summary}\centering\small
  \begin{tabular}{lcc@{\hspace{2em}}lcc}
    \toprule[1.5pt]
    指标 & 问题三 & 问题四 & 指标 & 问题三 & 问题四 \\
    \midrule[1pt]
    有效局数 & $50$ & $70$ & 平均 s/源（案例口径） & $252.0$ & $511.2$ \\
    清除源总数 & $667$ & $887$ & 平均 s/源（合并口径） & $247.5$ & $500.2$ \\
    清除比例 & $100\%$ & $100\%$ & 标准差／中位数 & $34.3$／$251.5$ & $83.7$／$515.8$ \\
    整局总时间均值 & $3301$ & $6338$ & 四分位区间 & $226.6\sim279.4$ & $445.7\sim563.9$ \\
    \quad 标准差／变异系数 & $291$／$8.8\%$ & $398$／$6.3\%$ & $\mathrm{P}_{95}$ & $301.6$ & $644.6$ \\
    动作数 & $127\sim201$ & $187\sim390$ & 最好／最差 & $185.2$／$314.8$ & $340.6$／$681.7$ \\
    程序运行时间/s & $0.4\sim2.2$ & $0.7\sim1.1$ & 清满 $16$ 源提前结束 & $5$ 局 & $3$ 局 \\
    \bottomrule[1.5pt]
  \end{tabular}
\end{table}

\begin{table}[!htbp]
  \caption{演练成绩按源数分层（平均定位清除时间，案例口径，虚拟秒/源）}\label{tab:practice-strat}\centering
  \begin{tabular}{lccccccc}
    \toprule[1.5pt]
    源数 $n$ & $10$ & $11$ & $12$ & $13$ & $14$ & $15$ & $16$ \\
    \midrule[1pt]
    问题三：局数 & $3$ & $12$ & $5$ & $7$ & $4$ & $6$ & $13$ \\
    \qquad\ \ 均值 & $296.4$ & $282.1$ & $277.9$ & $259.5$ & $247.5$ & $215.2$ & $218.3$ \\
    问题四：局数 & $10$ & $11$ & $13$ & $12$ & $10$ & $10$ & $4$ \\
    \qquad\ \ 均值 & $639.3$ & $575.5$ & $537.1$ & $494.7$ & $445.3$ & $418.8$ & $375.2$ \\
    \bottomrule[1.5pt]
  \end{tabular}
\end{table}

两张表合起来印证了\cref{sec:analysis:fixed}的判断。整局虚拟总时间的变异系数只有 $6\%\sim9\%$，而平均定位清除时间在 $n=10$ 与 $n=16$ 之间相差 $78\msec/$源（第三问）与 $264\msec/$源（第四问）——\textbf{波动几乎全部来自分母，而不是算法本身不稳定}。第四问 $70$ 局中有 $3$ 局在访问完 $19$ 或 $21$ 站时就清满 $16$ 个源而提前退出，这也是其总时间随源数轻微\emph{下降}的原因。

\subsection{本地模型与官方模拟器的一致性}

本文所有版本比较都在本地合成模拟器上完成，因此本地环境与官方环境的差距直接关系到这些结论能否外推。把\cref{tab:q4holdout}脚注中的本地留出分层预测与官方演练的同源数结果对照，得到\cref{tab:q4match}。

\begin{table}[!htbp]
  \caption{问题四本地留出预测与官方演练的分层对照（虚拟秒/源）}\label{tab:q4match}\centering
  \begin{tabular}{ccccc}
    \toprule[1.5pt]
    源数 $n$ & 本地留出预测 & 官方局数 & 官方实测 & 官方比本地 \\
    \midrule[1pt]
    $10$ & $642$ & $10$ & $639.3$ & $-0.4\%$ \\
    $13$ & $475$ & $12$ & $494.7$ & $+4.1\%$ \\
    $16$ & $346$ & $4$  & $375.2$ & $+8.4\%$ \\
    \bottomrule[1.5pt]
  \end{tabular}
\end{table}

两者在量级上吻合，但存在一个系统性趋势：\textbf{源数越多，官方环境比本地慢得越明显}，偏差从 $-0.4\%$ 升到 $+8.4\%$。一个合理的解释是官方案例中源的空间分布比本地采用的面积均匀分布更分散，使多目标之间顺路共享观测的机会变少；不过 $n=16$ 只有 $4$ 局，这一趋势本身也可能只是抽样噪声。

因此本文对这一比较只作如下有限的表述：\emph{在已有 $70$ 局官方演练覆盖的源数范围内，本地模型与官方模拟器给出的耗时量级一致，本地配对实验中“甲版本快于乙版本”这类定性结论可以谨慎外推。}它既不足以证明两种案例的分布相同，也不保证正式测试中同样的相对收益一定成立。

\begin{figure}[!htbp]
  \centering

  %% 画法建议：三个子图。(a) 散点图：横轴源数 $n$、纵轴整局虚拟总时间，第三问与第四问分别用不同颜色，
  %% 叠加回归直线 $T_{Q3}=2396+68n$、$T_{Q4}=6848-40n$ 及其 $95\%$ 带（数据取自 \texttt{q3\_runs.csv}、\texttt{q4\_runs.csv}，共 $50+70$ 点），
  %% 直观展示“总时间主要由固定成本决定”；第四问的回归线近乎水平甚至微降，可加注 $r=-0.18$。
  %% (b) 同样的横轴、纵轴改为平均定位清除时间 $\bar T$，画出 $120$ 个散点、按源数分组的箱线图，
  %% 并叠加双曲线 $\bar T=T_0/n+k$，说明分母效应。
  %% (c) 堆叠柱状图：第三问与第四问的每源时间构成（移动／检测／切频／光学定位／激光），数据取自\cref{tab:breakdown}。
  \fbox{\parbox[c][2.2cm][c]{0.86\textwidth}{\centering\small
  【待插图】（绘图要求见本段 \TeX{} 源码注释）}}
  \caption{演练成绩随源数的变化与时间构成}\label{fig:results}
\end{figure}

\cref{fig:results}把两问 $120$ 局演练成绩按源数与时间构成做了汇总。

\subsection{正式测试结果}

三次正式测试的结果分别按题目表 1 的格式列于\cref{tab:formal3}与\cref{tab:formal4}。

%%%%%%%%%%%%%%%%%%%%%%%%%%%%%%%%%%%%%%%%%%%%%%%%%%%%%%%%%%%%%%%%%%%%%%%%%%%%%%%%
%% TODO（提交前必改）：以下两张表在完成三次正式测试后按模拟器界面填写。
%%   1. 测试案例编码：直接抄模拟器界面显示的编码，不要改写。
%%   2. 清除干扰源个数：取 practice_logs/<局>/summary.json 的 cleared_count。
%%   3. 平均定位清除时间：取 average_per_cleared_s（= total_virtual_s / cleared_count）。
%%   4. 程序运行时间：取 program_elapsed_s（/enter 到测试结束的现实时长）。
%%   5. 填完后删除下面这一段 \noindent\fbox{...} 提示文字。
%%   6. 导出三次正式测试的加密日志（不要改文件名）放入支撑材料。
%%%%%%%%%%%%%%%%%%%%%%%%%%%%%%%%%%%%%%%%%%%%%%%%%%%%%%%%%%%%%%%%%%%%%%%%%%%%%%%%
\noindent\fbox{\parbox{\dimexpr\linewidth-14pt\relax}{\small\textbf{【待填写】}三次正式测试完成后，按模拟器界面的测试案例编码与本地
\texttt{summary.json} 中的 \texttt{cleared\_count}、\texttt{average\_per\_cleared\_s}、\texttt{program\_elapsed\_s}
填入下列两表，并删除本提示框。}}\medskip

\begin{table}[!htbp]
  \caption{问题三正式测试结果}\label{tab:formal3}\centering
  \begin{tabular}{cccc}
    \toprule[1.5pt]
    测试案例编码 & 清除干扰源个数 & 平均定位清除时间/s & 程序运行时间/s \\
    \midrule[1pt]
    测试 $1$ 案例编码 &  &  &  \\
    测试 $2$ 案例编码 &  &  &  \\
    测试 $3$ 案例编码 &  &  &  \\
    \bottomrule[1.5pt]
  \end{tabular}
\end{table}

\begin{table}[!htbp]
  \caption{问题四正式测试结果}\label{tab:formal4}\centering
  \begin{tabular}{cccc}
    \toprule[1.5pt]
    测试案例编码 & 清除干扰源个数 & 平均定位清除时间/s & 程序运行时间/s \\
    \midrule[1pt]
    测试 $1$ 案例编码 &  &  &  \\
    测试 $2$ 案例编码 &  &  &  \\
    测试 $3$ 案例编码 &  &  &  \\
    \bottomrule[1.5pt]
  \end{tabular}
\end{table}

\subsection{成绩的正确解读}

三次正式测试是三个独立随机案例，样本量极小，因此有必要先说清楚成绩该怎么读。\textbf{第一，三次成绩的差异主要由抽到的源数决定，而不是算法是否稳定}：演练中整局虚拟总时间的变异系数只有 $6\%\sim9\%$，而平均定位清除时间在 $n=10$ 与 $n=16$ 之间相差近一半（\cref{tab:practice-strat}），因此本文同时报告两个指标，解读时优先看整局总时间。\textbf{第二，两种平均口径不能混用}：案例口径与合并口径的定义见\cref{eq:twomeans}，两者相差约 $2\%$，本文所有对比都在同一口径下进行并在表注中标明。\textbf{第三，清除比例排在时间之前}：只有在 $N_c=n$ 时比较 $\bar T$ 才有意义，某局若未清完，$\bar T$ 反而可能因分母变小而“变好”，这类记录必须单列。\textbf{第四，程序运行时间不是效率指标}：它在演练中只有 $0.4\sim2.2\msec$（第三问）与 $0.7\sim1.1\msec$（第四问），相对 $1200\msec$ 的限制有三个数量级余量，主要由 HTTP 往返次数决定，与虚拟时间没有可比性。
